\begin{frame}{Correction encodage extraits \textit{Le crime de l'Orient-express}}
    
\end{frame}

\begin{frame}{Utiliser XPath: deux solutions}
    \begin{itemize}
        \item \href{https://www.oxygenxml.com/}{Oxygen}: solution payante pour faire de l'encodage en XML TEI, XPath et XSLT
        \item Extensions VS Code : \textit{XPath Notebook for Visual Studio Code} et \textit{XSLT/XPath for Visual Studio Code}
    \end{itemize}
\end{frame}

\begin{frame}{Installation Oxygen}
    \href{https://www.oxygenxml.com/xml_editor/download_oxygenxml_editor.html?os=Linux}{Télécharger Oxygen} puis faire une demande pour un essai gratuit de 30 jours.
\end{frame}
\begin{frame}{Installation pour utiliser XPath et XSLT dans VSCode}
    \begin{itemize}
        \item Installer Node.js: \url{https://nodejs.org/en/download}
        \item Dans VSCode, aller dans les extensions et installer:\\
        XPath Notebook for Visual Studio Code\\
        XSLT/XPath for Visual Studio Code
        \item Pour ouvrir un Note Book Xpath: ouvrir le fichier XML puis faire: ctrl+shift+p ou shift + pomme + p (pour mac).
    \end{itemize}
\end{frame}

\begin{frame}{XPath: définition}
\begin{itemize}
    \item \textbf{XPath =} \textit{XML Path Language}
    \item Langage de requête pour parcourir l'arbre XML
    \item Repose sur le principe de \textbf{\textcolor{red}{chemin}}
\end{itemize}
    
\end{frame}

\begin{frame}{Exemple}
\centering
\includegraphics[width=.9\linewidth]{03 - XPath/poeme_tei.png}

\end{frame}

\begin{frame}{Exemple}
    \begin{itemize}
        \item Quel est le titre du poème?\\
        \textcolor{blue}{\texttt{//title/text()}}
        \item Combien de strophes?\\
        \textcolor{blue}{\texttt{count(//lg)}}
        \item Que contient le deuxième vers de la deuxième strophe?
        \textcolor{blue}{\texttt{//lg[position()=2]/descendant::l [position()=2]/text()}}
    \end{itemize}
\end{frame}

\begin{frame}{Pourquoi utiliser XPath?}
\begin{itemize}
    \item \textbf{Navigation efficace }: XPath permet d'accéder rapidement à des éléments spécifiques d'un document XML sans avoir à parcourir toute la structure.
\item \textbf{Filtrage précis} : Avec XPath, on peut cibler exactement ce que l'on veut, que ce soit des titres, des chapitres ou des paragraphes.
\item \textbf{Simplicité }: Même si XML peut être complexe, XPath rend l'extraction d'informations beaucoup plus simple.
\end{itemize}

\end{frame}

